\chapter{Breast Reduction}

\section*{Introduction \& Clinical Indications}
Breast reduction, formally known as reduction mammoplasty, is a surgical procedure designed to alleviate the physical and psychological burdens associated with macromastia, or excessively large breasts. This operation involves the precise removal of excess glandular tissue, adipose tissue, and skin from the breasts, followed by reshaping and repositioning of the remaining breast mound and nipple-areola complex. The primary anatomical site of intervention is the breast itself, with careful preservation of the underlying chest wall musculature and neurovascular structures. Clinically, it is indicated for individuals experiencing chronic musculoskeletal pain, skin irritation, functional limitations, and significant psychological distress due to breast hypertrophy. As a definitive surgical intervention, breast reduction offers a durable solution for symptoms that have proven refractory to conservative management strategies, significantly improving patients' quality of life and physical comfort.

\begin{itemize}
  \item Reduction mammoplasty
  \item Reduction mammaplasty
  \item Therapeutic mammaplasty
  \item Gigantomastia correction
  \item Breast hypertrophy surgery
  \item Wise-pattern breast reduction
  \item Vertical (lollipop) reduction mammoplasty
\end{itemize}

The fundamental principle of breast reduction surgery revolves around the concept of excising excess breast tissue and skin while preserving a robust blood supply and nerve innervation to the nipple-areola complex (NAC). This is achieved by designing a pedicle, a segment of breast tissue containing the NAC, which remains attached to the chest wall and carries its vital neurovascular supply. Common pedicle designs include the inferior pedicle, superior pedicle, and superomedial pedicle, each chosen based on breast size, desired reduction volume, and surgeon preference. The rationale is to reduce the overall weight and volume of the breast, elevate the NAC to a more aesthetic and functional position, and reshape the remaining breast tissue into a smaller, more conical, and proportionate mound. Technical goals include achieving bilateral symmetry, maintaining NAC viability and sensation, minimizing visible scarring, and creating a stable, long-lasting aesthetic result that alleviates the patient's symptoms. The careful planning of skin incisions and tissue resection patterns, such as the inverted-T (Wise pattern) or vertical (lollipop) pattern, is critical to achieving these objectives while placing scars in the least conspicuous locations.

The history of breast reduction surgery dates back to the late 19th and early 20th centuries, with early attempts focused primarily on tumor removal rather than aesthetic reduction. One of the earliest documented cases of breast reduction for macromastia was performed by Thorek in 1922, involving a large wedge resection. Significant advancements began in the 1920s and 1930s with surgeons like Biesenberger (1928) who described a superior pedicle technique, and Schwarzmann (1930) who emphasized the importance of preserving the NAC's blood supply.

The mid-20th century saw the development of more sophisticated pedicle-based techniques. Strombeck (1960) introduced the bipedicle flap, which significantly improved NAC viability. Pitanguy (1960s) popularized a keyhole pattern with a superior pedicle, offering a more aesthetic outcome. The 1970s marked a major turning point with the widespread adoption of the inferior pedicle technique, notably refined by Robbins and others, which became a cornerstone of modern reduction mammoplasty due to its reliability, especially for larger reductions. The 1990s brought the advent of vertical reduction mammoplasty, championed by Lejour, which aimed to reduce the horizontal scar in the inframammary fold, leading to shorter scars and often a more youthful breast shape. Contemporary techniques continue to evolve, incorporating liposuction as an adjunct for contouring and exploring minimal-scar approaches, constantly striving for improved safety, aesthetics, and patient satisfaction.

Breast reduction surgery is indicated for individuals experiencing significant physical and psychological symptoms directly attributable to macromastia. These indications include:

\begin{itemize}
  \item Chronic neck, upper back, and shoulder pain that is refractory to conservative management, often exacerbated by the weight of the breasts.
  \item Deep shoulder grooving or indentation caused by tight bra straps, leading to nerve compression or skin irritation.
  \item Recurrent intertrigo, rashes, or fungal infections in the inframammary fold due to skin-on-skin contact and moisture retention.
  \item Significant functional limitations, including difficulty participating in physical activities, exercise, or sports.
  \item Impairment with daily hygiene, finding properly fitting clothing, and maintaining an active lifestyle.
  \item Peripheral nerve symptoms such as numbness, tingling, or paresthesias in the hands and arms, occasionally associated with severe breast hypertrophy.
  \item Psychological distress, body image dissatisfaction, and social anxiety related to breast size and shape.
  \item Asymmetry between breasts that causes significant physical discomfort or psychological distress.
\end{itemize}

Breast reduction is primarily positioned as a first-line surgical treatment for symptomatic macromastia, rather than a secondary or salvage procedure. It is the definitive operative solution chosen when a patient's symptoms are severe, chronic, and have not responded adequately to non-surgical interventions such as supportive brassieres, physical therapy, weight management, or topical treatments for skin irritation. In this context, it directly addresses the underlying anatomical cause of the symptoms by reducing breast volume and weight.

While its main role is primary, breast reduction can occasionally serve a secondary or adjunct function. For instance, it may be performed to correct significant breast asymmetry that causes functional issues or psychological distress, either as a standalone procedure or following a previous breast operation. It can also be considered as a balancing procedure in breast reconstruction, where one breast is reduced to match the size and shape of a reconstructed contralateral breast. However, its most common and impactful role remains as the initial and definitive surgical remedy for the debilitating effects of oversized breasts.

\section*{Relevant Surgical Anatomy}
The breast is a modified apocrine gland situated on the anterior chest wall, extending from the second to the sixth rib vertically and from the sternal edge to the midaxillary line horizontally. It comprises glandular tissue (lobules and ducts), adipose tissue, and fibrous connective tissue (Cooper's ligaments) encased within a skin envelope. The nipple-areola complex (NAC) is centrally located, containing smooth muscle fibers and sensory nerve endings.

The arterial supply to the breast is rich and originates from several sources:
\begin{itemize}
  \item Internal mammary artery perforators (medial aspect).
  \item Lateral thoracic artery (lateral aspect).
  \item Thoracoacromial artery (superior aspect).
  \item Intercostal artery perforators (inferior aspect).
\end{itemize}
Venous drainage largely parallels the arterial supply, emptying into the axillary and internal mammary veins. Lymphatic drainage is predominantly to the axillary lymph nodes (75\%), with the remainder draining to the internal mammary nodes and supraclavicular nodes. Sensory innervation to the breast and NAC is primarily provided by the anterior and lateral cutaneous branches of the second to sixth intercostal nerves, with the fourth intercostal nerve being crucial for nipple sensation. Surgical landmarks include the sternal notch, midclavicular line, inframammary fold (IMF), and the desired position of the new nipple. Understanding these intricate anatomical relationships is paramount for preserving NAC viability and sensation during reduction mammoplasty.

\section*{Preoperative Workup \& Preparation}
A comprehensive preoperative workup is essential to ensure patient safety and optimize outcomes for breast reduction surgery. This typically includes:

\begin{itemize}
  \item \textbf{Medical Evaluation:} A thorough history and physical examination, including a detailed breast examination, is performed. Patients undergo a medical clearance from their primary care physician, and often a cardiac and anesthesia assessment, especially for those with comorbidities.
  \item \textbf{Laboratory Tests:} Routine blood tests, such as a complete blood count, basic metabolic panel, and coagulation studies, are ordered.
  \item \textbf{Breast Imaging:} Mammography is typically required for women over 40 or those with a family history of breast cancer, and sometimes for younger patients with specific concerns, to rule out any underlying breast pathology. Ultrasound may also be used.
  \item \textbf{Medication Review:} Patients are instructed to discontinue medications that increase bleeding risk (e.g., aspirin, NSAIDs, warfarin, certain herbal supplements) for at least two weeks prior to surgery.
  \item \textbf{Smoking Cessation:} Strict smoking cessation is mandated for at least 4-6 weeks preoperatively, as smoking significantly impairs wound healing and increases the risk of nipple necrosis and other complications.
  \item \textbf{NPO Status:} Patients must adhere to strict nil per os (NPO) guidelines for food and drink for a specified period (typically 6-8 hours) before general anesthesia.
  \item \textbf{Prophylactic Antibiotics:} Intravenous prophylactic antibiotics are administered shortly before the incision to reduce the risk of surgical site infection.
  \item \textbf{Informed Consent:} A detailed discussion covers the surgical plan, expected outcomes, potential risks (including permanent scars, changes in sensation, and impact on breastfeeding), and the recovery process. Written informed consent is obtained.
  \item \textbf{Logistical Planning:} Patients are advised to arrange for assistance at home during the initial recovery period and to plan for time off work or school.
\end{itemize}

Careful consideration of contraindications and precautions is vital for patient safety and optimal surgical outcomes.

\textbf{Absolute Contraindications:}
\begin{itemize}
  \item Active breast infection or inflammatory breast conditions.
  \item Uncontrolled severe medical comorbidities (e.g., unstable angina, recent myocardial infarction, severe uncontrolled diabetes, severe pulmonary disease) that render general anesthesia or surgery excessively risky.
  \item Pregnancy or active breastfeeding.
  \item Unrealistic patient expectations regarding scar appearance, breast size, or shape.
  \item Unwillingness to accept the permanent scars inherent to the procedure.
  \item Undiagnosed breast lump or suspicious mammographic findings that require further investigation.
\end{itemize}

\textbf{Relative Contraindications \& Precautions:}
\begin{itemize}
  \item \textbf{Active Smoking:} While not an absolute contraindication, active smoking significantly increases the risk of wound healing complications, nipple necrosis, and infection. Patients are strongly advised to cease smoking for at least 4-6 weeks preoperatively.
  \item \textbf{Obesity (BMI > 30-35 kg/m\textasciicircum2):} Increased BMI is associated with higher rates of complications such as infection, wound dehiscence, fat necrosis, and venous thromboembolism. Weight loss may be recommended prior to surgery.
  \item \textbf{Poorly Controlled Diabetes Mellitus:} Increases infection risk and impairs wound healing. Strict glycemic control is essential.
  \item \textbf{Significant Personal or Family History of Breast Cancer:} Requires thorough preoperative screening and discussion with the patient regarding future breast cancer surveillance, as extensive tissue rearrangement can alter mammographic patterns.
  \item \textbf{Desire for Future Breastfeeding:} While not always precluded, breast reduction can compromise the ability to breastfeed. This must be thoroughly discussed with patients of childbearing age.
  \item \textbf{Psychological Instability or Body Dysmorphic Disorder:} Patients with significant psychological issues may not be appropriate candidates and require psychiatric evaluation.
  \item \textbf{Extremely Large Reductions:} Very large reductions (e.g., >1000g per breast) may necessitate a free nipple graft technique, which results in complete loss of nipple sensation and breastfeeding ability.
\end{itemize}

\section*{Surgical Technique \& Procedure}
Breast reduction surgery is typically performed under general anesthesia, with the patient positioned supine and arms abducted on arm boards to allow full access to the chest and axillae.

The procedure generally follows these key steps:
\begin{enumerate}
  \item \textbf{Preoperative Markings:} The surgeon meticulously marks the planned incision lines, the new nipple position, and the areas of tissue resection while the patient is upright. Common patterns include the inverted-T (Wise pattern), which involves circumareolar, vertical, and inframammary fold incisions, or the vertical (lollipop) pattern, which omits the horizontal inframammary incision.
  \item \textbf{Incision and Skin Flap Elevation:} Incisions are made through the skin and subcutaneous tissue according to the markings. Skin flaps are carefully elevated to expose the underlying glandular and adipose breast tissue.
  \item \textbf{Tissue Resection and Pedicle Creation:} Excess glandular and adipose tissue is precisely resected from the periphery of the breast. A central pedicle, containing the nipple-areola complex (NAC) and its neurovascular supply, is carefully preserved and sculpted. The specific pedicle design (e.g., inferior, superior, superomedial) is chosen based on the degree of reduction and surgeon's preference to ensure optimal NAC viability.
  \item \textbf{NAC Repositioning:} The NAC, still attached to its pedicle, is transposed to its new, higher position on the breast mound, typically marked preoperatively. In cases of extreme macromastia, a free nipple graft technique may be employed, where the NAC is completely detached and reattached as a skin graft.
  \item \textbf{Breast Reshaping and Contouring:} The remaining breast tissue is meticulously reshaped and tightened to create a smaller, more conical, and aesthetically pleasing breast mound. This often involves suturing the glandular tissue to the pectoralis fascia for support and projection.
  \item \textbf{Hemostasis:} Meticulous hemostasis is achieved throughout the procedure using electrocautery to minimize bleeding and reduce the risk of hematoma formation.
  \item \textbf{Wound Closure:} The skin flaps are redraped, and the incisions are closed in multiple layers using absorbable sutures for the deeper tissues and often non-absorbable or absorbable sutures for the skin. The NAC is sutured into its new position.
  \item \textbf{Dressing and Drains:} A sterile dressing and a supportive surgical bra are applied. Surgical drains (e.g., Jackson-Pratt) may be placed in the inframammary fold to collect serosanguinous fluid, though many surgeons opt to omit them in selected cases.
\end{enumerate}
The entire procedure typically lasts between two and four hours, depending on the complexity and volume of tissue to be removed.

\section*{Complications \& Risks}
As with any surgical procedure, breast reduction carries potential complications and risks, which can be broadly categorized into general surgical risks and procedure-specific risks.

\textbf{General Surgical Complications:}
\begin{itemize}
  \item \textbf{Anesthesia Risks:} Reactions to anesthetic agents, respiratory or cardiac complications.
  \item \textbf{Bleeding/Hematoma:} Accumulation of blood under the skin, potentially requiring reoperation.
  \item \textbf{Infection:} Surgical site infection, requiring antibiotics or drainage.
  \item \textbf{Deep Vein Thrombosis (DVT) / Pulmonary Embolism (PE):} Blood clots in the legs that can travel to the lungs, a rare but serious complication.
\item \textbf{Seroma:} Accumulation of clear fluid under the skin, sometimes requiring aspiration.
\end{itemize}

\textbf{Procedure-Specific Complications:}
\begin{itemize}
  \item \textbf{Scarring:} Permanent scars are inevitable. These can be prominent, hypertrophic (raised and red), or keloid (excessively overgrown), especially in individuals with a predisposition.
  \item \textbf{Changes in Nipple/Breast Sensation:} Partial or complete loss of sensation in the nipple-areola complex or breast skin is common and can be permanent. Conversely, some patients may experience hypersensitivity.
  \item \textbf{Fat Necrosis:} Areas of hardened, non-viable fat tissue within the breast, which can cause lumps, pain, and sometimes require excision.
  \item \textbf{Wound Healing Problems/Dehiscence:} Delayed healing or separation of incision lines, particularly at the T-junctions of the inverted-T pattern, more common in smokers or diabetics.
  \item \textbf{Asymmetry:} While efforts are made to achieve symmetry, some degree of asymmetry in breast size, shape, or nipple position can persist or develop.
  \item \textbf{Partial or Complete Nipple-Areola Complex (NAC) Necrosis:} A rare but devastating complication where the blood supply to the NAC is compromised, leading to tissue death. This is more common with very large reductions or in smokers.
  \item \textbf{Inability to Breastfeed:} Breast reduction surgery can damage milk ducts and nerves, potentially impairing or preventing future breastfeeding.
  \item \textbf{Contour Irregularities/Deformities:} Unevenness or depressions in the breast contour.
  \item \textbf{Skin Discoloration:} Bruising and temporary or permanent changes in skin pigmentation.
\end{itemize}

\section*{Postoperative Care \& Recovery}
The postoperative recovery from breast reduction surgery typically follows a predictable timeline, though individual experiences may vary.

\textbf{Immediate Postoperative Period (PACU to 48 hours):}
Upon waking from general anesthesia, patients are monitored in the Post-Anesthesia Care Unit (PACU) for vital signs, pain, and nausea. A surgical bra or compression dressing is applied. Pain is managed with prescribed oral analgesics, and antiemetics are given as needed. Drains, if placed, are typically monitored for output. Patients are encouraged to ambulate gently to reduce the risk of deep vein thrombosis. Most patients are discharged home the same day or after an overnight hospital stay. Swelling and bruising are expected, and a feeling of tightness in the chest is common.

\textbf{Early Recovery (First 1-2 Weeks):}
During this period, patients will experience continued swelling, bruising, and tenderness. Pain gradually subsides and can usually be managed with over-the-counter pain relievers by the end of the first week. Activity is significantly restricted; patients are advised to avoid lifting arms above shoulder level, lifting heavy objects (>5-10 lbs), strenuous exercise, and sleeping on their stomach. The surgical bra must be worn continuously, often for 4-6 weeks. Drains, if present, are usually removed within the first few days to a week. Incision lines are kept clean and dry, and patients are typically advised to avoid submerging the breasts in baths or pools. Return to light, non-strenuous work is often possible within 1-2 weeks.

\textbf{Mid-Term Recovery (2-6 Weeks):}
Swelling continues to resolve, and the breasts begin to soften and settle into their new shape. Scars will appear red and raised but will gradually mature. Most patients can gradually resume more normal activities, including light exercise, after 3-4 weeks, with full strenuous exercise typically permitted after 6 weeks, provided healing is progressing well. Scar management, such as massage or silicone sheeting, may be initiated around this time.

\textbf{Long-Term Recovery (3-6 Months and Beyond):}
The final breast shape and size will become apparent as all swelling resolves, which can take up to 6 months or even a year. Scars will continue to fade and flatten, though they are permanent. Nipple sensation may gradually return over several months, but some degree of permanent numbness is common. Patients can typically resume all normal activities without restriction. Regular follow-up appointments are crucial to monitor healing and address any concerns.

During breast reduction surgery, the surgeon meticulously documents several key findings. This includes the weight of the resected tissue from each breast, which is crucial for medical necessity documentation and for assessing the degree of reduction achieved. The surgeon also notes the pedicle type used for the nipple-areola complex (NAC), the estimated blood supply to the NAC, and the overall symmetry achieved intraoperatively. Any unexpected findings, such as suspicious lesions or unusual tissue characteristics, are also recorded.

All resected breast tissue is routinely sent for histopathological examination. The pathology report typically confirms the presence of benign glandular and adipose breast tissue, often describing features consistent with macromastia, such as stromal hypertrophy and lobular hyperplasia. The absence of malignant or atypical cells is reassuring for both the patient and the surgeon. In cases where a suspicious lesion was noted preoperatively or intraoperatively, the pathologist will provide a definitive diagnosis, guiding any further management if necessary. This pathological confirmation is a standard and essential component of the procedure, ensuring that no underlying breast pathology is overlooked.

A normal and successful postoperative course following breast reduction is characterized by a progressive reduction in pain, swelling, and bruising, leading to the gradual emergence of the desired breast size and shape. Patients can expect initial discomfort to be well-controlled with prescribed oral analgesics, with pain steadily diminishing over the first week. Swelling and bruising will be most prominent in the first few days and will gradually subside over several weeks. The incision lines are expected to heal cleanly, without signs of infection, excessive redness, or dehiscence. The breasts will initially feel firm and tight, but will soften and settle into a more natural contour over the first few months. A key indicator of a normal course is the preservation of nipple-areola complex viability, with good color and capillary refill. Patients should experience significant relief from their preoperative symptoms of neck, back, and shoulder pain, as well as improved comfort and ability to participate in physical activities, marking a successful functional outcome.

Comprehensive postoperative care and diligent follow-up are crucial for optimal healing and long-term satisfaction after breast reduction.

\begin{itemize}
  \item \textbf{Follow-up Visits:} Patients typically have their first follow-up visit within 1-2 weeks post-surgery to check wound healing, remove any drains, and address immediate concerns. Subsequent visits are scheduled at 1 month, 3 months, 6 months, and 1 year to monitor scar maturation, assess breast shape, and evaluate long-term outcomes.
  \item \textbf{Suture/Staple Removal:} If non-absorbable sutures or staples were used, they are typically removed at the first or second follow-up visit.
  \item \textbf{Wound Care:} Patients are instructed on how to care for their incision sites, which usually involves keeping them clean and dry. Specific instructions regarding showering and bathing will be provided.
  \item \textbf{Supportive Bra:} A surgical support bra is worn continuously for 4-6 weeks to minimize swelling and provide support to the healing tissues.
  \item \textbf{Activity Restrictions:} Gradual return to normal activities is advised. Strenuous exercise, heavy lifting, and activities that raise the heart rate or blood pressure are restricted for at least 4-6 weeks.
  \item \textbf{Scar Management:} Once incisions are fully healed, typically around 3-4 weeks, scar massage, silicone sheeting, or topical scar creams may be recommended to improve scar appearance.
  \item \textbf{Warning Signs:} Patients are educated on warning signs that necessitate immediate contact with the surgeon, including fever, increasing pain not relieved by medication, excessive redness or warmth around the incisions, pus-like drainage, significant swelling or bruising, or any changes in nipple color or sensation.
  \item \textbf{Long-term Surveillance:} Routine breast cancer screening, including mammography, should continue according to age and risk factors, with the understanding that prior surgery may alter mammographic interpretation.
\end{itemize}

\section*{Outcomes \& Prognosis}
Symptomatic macromastia is a prevalent condition globally, affecting a significant proportion of women across various age groups, though it is most commonly observed in women of childbearing age and beyond. The exact incidence is difficult to quantify precisely due to varying definitions of macromastia, but surveys indicate that a substantial number of women experience physical and psychological symptoms related to large breasts. Studies consistently report high frequencies of chronic neck, back, and shoulder pain (up to 80-90\%), shoulder grooving (60-70\%), and inframammary skin irritation (40-50\%) among affected individuals. Breast reduction surgery is one of the most frequently performed reconstructive and aesthetic plastic surgery procedures worldwide, reflecting the high demand for relief from these symptoms. While the procedure is predominantly performed on women, a small percentage of men with gynecomastia also undergo similar reduction techniques. Morbidity rates are generally low, with major complications occurring in a small percentage of cases, typically higher in patients with significant medical comorbidities such as obesity, diabetes, and especially in active smokers. Mortality from breast reduction surgery is exceedingly rare, estimated to be less than 0.01%.

Breast reduction surgery consistently yields exceptionally high rates of patient satisfaction and significant improvement in quality of life. Numerous studies and patient surveys report satisfaction rates exceeding 90-95\%. The primary functional outcomes, such as relief from chronic neck, back, and shoulder pain, resolution of shoulder grooving, and improvement in inframammary skin irritation, are typically achieved in the vast majority of patients. Functional limitations related to physical activity and clothing fit also resolve, leading to enhanced participation in daily life and exercise.

The cosmetic result is generally durable over the long term, with breasts maintaining a smaller, more proportionate shape. While some mild ptosis (drooping) may naturally occur with aging, gravity, and weight fluctuations, the significant reduction in volume is permanent. Prognosis for symptomatic relief and improved quality of life is excellent. Recurrence of macromastia is uncommon but can occur with substantial weight gain or hormonal changes. Key prognostic factors for optimal outcomes include patient selection, adherence to preoperative and postoperative instructions (especially smoking cessation), and the skill of the surgical team. While specific landmark clinical trials like those for appendicitis (e.g., APPAC) are not directly applicable to breast reduction, large cohort studies and systematic reviews consistently demonstrate the procedure's efficacy and positive impact on patient-reported outcomes and health-related quality of life.

\section*{References}
\begin{enumerate}
  \item MedlinePlus. Breast reduction. U.S. National Library of Medicine; 2024. Available from: \url{https://medlineplus.gov/ency/article/002941.htm}
  \item Sabiston Textbook of Surgery: The Biological Basis of Modern Surgical Practice. 22nd ed. Elsevier; 2022.
  \item Grabb and Smith's Plastic Surgery. 8th ed. Wolters Kluwer; 2019.
  \item American Society of Plastic Surgeons. Clinical Practice Guideline for Breast Reduction. 2023. Available from: \url{https://www.plasticsurgery.org/for-medical-professionals/clinical-resources/clinical-practice-guidelines}
\end{enumerate}

\clearpage
